\chapter{THIẾT KẾ HỆ THỐNG PHẦN MỀM VÀ CƠ SỞ DỮ LIỆU}
\label{ch:chuong3}

\section{Kiến trúc phần mềm phân tầng Backend NestJS}
\label{sec:kien_truc_backend}

Hệ thống máy chủ backend của SmartDorm được xây dựng trên nền tảng NestJS Framework (ngôn ngữ TypeScript) theo mô hình kiến trúc phân tầng Clean Architecture kết hợp hệ thống mô-đun hóa (Modular Design). Kiến trúc này tách biệt triệt để giữa tầng định tuyến giao tiếp (Controllers/Gateways), tầng nghiệp vụ logic cốt lõi (Services) và tầng truy xuất cơ sở dữ liệu (Prisma ORM Repositories).

Hệ thống backend được chia thành các phân hệ mô-đun chức năng độc lập:
\begin{itemize}
    \item \textbf{AuthModule \& UsersModule:} Xử lý quy trình đăng nhập, mã hóa mật khẩu Bcrypt, cấp phát token JWT và thực thi phân quyền người dùng RBAC.
    \item \textbf{RoomsModule \& BuildingsModule:} Quản lý cơ cấu phòng ở, danh mục tòa nhà, sức chứa và liên kết phần cứng vi điều khiển (\texttt{espDeviceId}).
    \item \textbf{DevicesModule:} Hiện thực các bộ điều khiển thiết bị (Đèn, Quạt, Khóa, Còi), quản lý lịch sử lệnh và thực thi cơ chế Hoàn tác (Undo).
    \item \textbf{AccessLogModule:} Tiếp nhận dữ liệu quẹt thẻ từ phần cứng, đối soát quyền vào phòng và ghi nhận lịch sử ra vào an ninh.
    \item \textbf{EmergencyModule:} Vận hành bộ quy tắc an toàn PCCC (Rule Engine), tự động ngắt điện phòng khi có sự cố khí gas và quản lý quy trình khắc phục sự cố.
    \item \textbf{AnalyticsModule:} Tổng hợp các chỉ số báo cáo vi khí hậu, an ninh, sự cố và xuất báo cáo dữ liệu định dạng Excel/CSV.
    \item \textbf{IntelligenceModule:} Vận hành các mô hình học máy phát hiện bất thường môi trường (Isolation Forest) và trợ lý ảo thông minh AI Copilot.
    \item \textbf{GatewayModule \& MqttModule:} Đảm nhận truyền thông thời gian thực với thiết bị phần cứng qua MQTT và giao diện web qua WebSockets Socket.IO.
\end{itemize}

\section{Thiết kế Cơ sở Dữ liệu quan hệ với Prisma ORM}
\label{sec:thiet_ke_csdl}

Cơ sở dữ liệu của SmartDorm được thiết kế chuẩn hóa ở dạng chuẩn 3NF với 12 thực thể quan hệ, được quản lý chặt chẽ thông qua Prisma ORM. Thiết kế này vừa bảo đảm tính toàn vẹn dữ liệu, vừa tối ưu hóa tốc độ truy vấn đối với các bản ghi lịch sử chuỗi thời gian lớn.

Bảng \ref{tab:thuc_the_csdl} tóm tắt các thực thể dữ liệu quan trọng nhất trong hệ thống.

\begin{table}[H]
\centering
\small
\begin{tabularx}{\textwidth}{|c|p{2.6cm}|p{4.0cm}|X|}
\hline
\textbf{STT} & \textbf{Tên bảng (Entity)} & \textbf{Các trường khóa chính/ngoại} & \textbf{Ý nghĩa và chức năng lưu trữ} \\
\hline
1 & \texttt{users} & \texttt{id (PK), studentId, roomId (FK)} & Lưu thông tin sinh viên, ban quản lý, mã thẻ RFID, mật khẩu băm và vai trò RBAC \\
\hline
2 & \texttt{buildings} & \texttt{id (PK), name} & Quản lý danh mục các khối nhà KTX (Tòa A, Tòa B, Tòa C) và số tầng \\
\hline
3 & \texttt{rooms} & \texttt{id (PK), buildingId (FK), espDeviceId} & Lưu thông tin số phòng, tầng, sức chứa và mã định danh phần cứng ESP32 liên kết \\
\hline
4 & \texttt{devices} & \texttt{id (PK), roomId (FK), type} & Quản lý danh mục 4 thiết bị trong phòng (LIGHT, FAN, DOOR\_LOCK, BUZZER) \\
\hline
5 & \texttt{sensor\_data} & \texttt{id (PK), roomId (FK), createdAt} & Lưu trữ chuỗi thời gian nhiệt độ, độ ẩm, nồng độ khí gas và cờ cảnh báo an toàn \\
\hline
6 & \texttt{access\_logs} & \texttt{id (PK), userId (FK), roomId (FK)} & Lưu lịch sử quẹt thẻ RFID (UID thẻ, kết quả Granted/Denied, mốc thời gian) \\
\hline
7 & \texttt{emergencies} & \texttt{id (PK), roomId (FK), level} & Lưu trữ các sự cố PCCC (FIRE, GAS\_LEAK), nồng độ cảm biến và trạng thái xử lý \\
\hline
8 & \texttt{command\_logs} & \texttt{id (PK), deviceId (FK), userId (FK)} & Lưu vết toàn bộ lệnh điều khiển (bật/tắt/mở khóa) phục vụ cơ chế Hoàn tác Undo \\
\hline
\end{tabularx}
\caption{Danh mục các bảng thực thể dữ liệu cốt lõi trong hệ thống SmartDorm}
\label{tab:thuc_the_csdl}
\end{table}

\begin{figure}[H]
\centering
\begin{tikzpicture}[
    entity/.style={rectangle, draw=blue!80!black, fill=blue!5, thick, minimum width=3.8cm, minimum height=1.8cm, rounded corners=3pt, align=center, font=\scriptsize},
    rel/.style={-{Stealth[length=2.5mm]}, thick, color=blue!70!black}
]
    % Entities
    \node[entity] (bld) at (0, 6) {\textbf{BUILDING}\\\texttt{id (PK)}\\\texttt{name (UQ)}\\\texttt{floors}};
    \node[entity] (room) at (5.5, 6) {\textbf{ROOM}\\\texttt{id (PK)}\\\texttt{roomNumber}\\\texttt{buildingId (FK)}\\\texttt{espDeviceId (UQ)}};
    \node[entity] (user) at (11, 6) {\textbf{USER}\\\texttt{id (PK)}\\\texttt{studentId (UQ)}\\\texttt{rfidCardId (UQ)}\\\texttt{role, roomId (FK)}};
    
    \node[entity] (dev) at (0, 2) {\textbf{DEVICE}\\\texttt{id (PK)}\\\texttt{roomId (FK)}\\\texttt{type, isOn}\\\texttt{isOnline}};
    \node[entity] (sensor) at (5.5, 2) {\textbf{SENSOR\_DATA}\\\texttt{id (PK)}\\\texttt{roomId (FK)}\\\texttt{temperature, humidity}\\\texttt{gasLevel, gasAlert}};
    \node[entity] (emg) at (11, 2) {\textbf{EMERGENCY}\\\texttt{id (PK)}\\\texttt{roomId (FK)}\\\texttt{type, level}\\\texttt{sensorValue, isResolved}};

    \node[entity] (cmd) at (2.5, -1.8) {\textbf{COMMAND\_LOG}\\\texttt{id (PK)}\\\texttt{deviceId (FK)}\\\texttt{userId (FK)}\\\texttt{action, payload, status}};
    \node[entity] (acc) at (8.5, -1.8) {\textbf{ACCESS\_LOG}\\\texttt{id (PK)}\\\texttt{roomId (FK)}\\\texttt{userId (FK)}\\\texttt{rfidCardId, action}};

    % Relationships
    \draw[rel] (bld) -- node[above, font=\tiny] {1 - N} (room);
    \draw[rel] (room) -- node[above, font=\tiny] {1 - N} (user);
    \draw[rel] (room) -- node[left, font=\tiny] {1 - N} (dev);
    \draw[rel] (room) -- node[right, font=\tiny] {1 - N} (sensor);
    \draw[rel] (room) -- node[right, font=\tiny] {1 - N} (emg);
    \draw[rel] (dev) -- node[left, font=\tiny] {1 - N} (cmd);
    \draw[rel] (room) -- node[right, font=\tiny] {1 - N} (acc);
    \draw[rel] (user) -- node[right, font=\tiny] {1 - N} (acc);
\end{tikzpicture}
\caption{Sơ đồ mối quan hệ thực thể (Entity-Relationship Diagram) của cơ sở dữ liệu SmartDorm}
\label{fig:erd_diagram}
\end{figure}

\section{Các thuật toán và mẫu thiết kế phần mềm cốt lõi}
\label{sec:thuat_toan_patterns}

\subsection{Thuật toán sinh mã thẻ RFID tự động theo sinh viên}
\label{subsec:thuat_toan_rfid}

Trong thực tế triển khai hệ thống kiểm soát an ninh thẻ từ, việc nhập thủ công chuỗi UID vật lý của từng thẻ thường tốn kém thời gian và dễ xảy ra sai sót nhập liệu. Để tự động hóa quy trình này, SmartDorm hiện thực thuật toán sinh mã thẻ RFID tất định (Deterministic RFID Generation Algorithm).

\begin{itemize}
    \item \textbf{Cơ sở lý thuyết:} Thẻ Mifare RC522 1K có định danh vật lý UID gồm 4 bytes dữ liệu, tương ứng với chuỗi 8 ký tự Hex (hệ cơ số 16).
    \item \textbf{Nguyên tắc giải thuật:} Sử dụng hàm băm mật mã an toàn SHA-256 mã hóa chuỗi kết hợp giữa tiền tố hệ thống bí mật (\texttt{salt = "SMARTDORM\_RFID\_"}) và Mã số sinh viên (\texttt{studentId}). Sau đó, trích xuất 8 ký tự Hex đầu tiên viết hoa và ghép với tiền tố chuẩn \texttt{CARD\_}.
    \item \textbf{Đặc tính an toàn:} Giải thuật có tính tất định (cùng một mã sinh viên luôn tạo ra một mã thẻ duy nhất), không thể đảo ngược, không trùng lặp và tương thích hoàn toàn với các đầu đọc RFID chuẩn Mifare.
\end{itemize}

Đoạn mã TypeScript hiện thực thuật toán trong tệp \texttt{rfid-utils.ts} được thể hiện trong Listing \ref{lst:rfid_algorithm}:

\begin{lstlisting}[language=Java, caption={Thuật toán sinh mã thẻ RFID tự động dựa trên SHA-256}, label={lst:rfid_algorithm}]
import * as crypto from 'crypto';

export function generateStudentRfid(studentId: string): string {
  const cleanId = (studentId || '').trim();
  if (!cleanId) {
    const randomHex = crypto.randomBytes(4).toString('hex').toUpperCase();
    return `CARD_${randomHex}`;
  }
  // Bam SHA-256 chuoi salt ket hop MSSV
  const hash = crypto.createHash('sha256')
                     .update(`SMARTDORM_RFID_${cleanId}`)
                     .digest('hex');
  // Trich xuat 8 ky tu Hex dau tien tao ma the UID 4 bytes chuan
  const uidHex = hash.substring(0, 8).toUpperCase();
  return `CARD_${uidHex}`;
}
\end{lstlisting}

\subsection{Mẫu thiết kế Command Pattern và cơ chế Hoàn tác lệnh (Undo)}
\label{subsec:command_pattern}

Khi quản trị viên hoặc sinh viên thực hiện thao tác điều khiển thiết bị (Bật đèn, Tắt quạt, Mở khóa), việc gửi trực tiếp lệnh điều khiển xuống phần cứng mà không quản lý lịch sử sẽ dẫn đến rủi ro: người dùng bấm nhầm nhưng không thể khôi phục lại trạng thái trước đó, hoặc khi xảy ra sự cố không thể truy vết ai đã gửi lệnh.

Hệ thống áp dụng mẫu thiết kế \textbf{Command Pattern} kết hợp ngăn xếp hoàn tác (Undo Stack) với cấu trúc như sau:
\begin{enumerate}
    \item \textbf{Giao diện ICommand:} Định nghĩa hai phương thức cốt lõi là \texttt{execute()} và \texttt{undo()}, cùng phương thức \texttt{getDescription()} mô tả chi tiết thao tác.
    \item \textbf{Các Concrete Commands:} Các lớp lệnh cụ thể bao gồm \texttt{TurnOnLightCommand}, \texttt{TurnOffLightCommand}, \texttt{ToggleFanCommand}, \texttt{ControlDoorCommand}. Mỗi lớp lưu giữ trạng thái trước đó của thiết bị để phục vụ khôi phục.
    \item \textbf{CommandInvokerService:} Đóng vai trò là đối tượng điều phối (Invoker). Khi một lệnh được thực thi thành công, đối tượng lệnh sẽ được đẩy vào ngăn xếp lịch sử (\texttt{commandHistory}).
    \item \textbf{Cơ chế Undo:} Khi người dùng nhấn nút ``Hoàn tác'' trên giao diện web, hệ thống lấy lệnh gần nhất ra khỏi ngăn xếp (\texttt{pop()}) và kích hoạt phương thức \texttt{undo()} tương ứng, gửi lệnh MQTT đảo ngược trạng thái thiết bị và cập nhật đồng bộ lên cơ sở dữ liệu.
\end{enumerate}

\begin{lstlisting}[language=Java, caption={Hiện thực Command Invoker và cơ chế Undo trong NestJS}, label={lst:command_invoker}]
@Injectable()
export class CommandInvokerService {
  private commandHistory: ICommand[] = [];
  private readonly MAX_HISTORY = 50;

  async executeCommand(command: ICommand): Promise<CommandResult> {
    const result = await command.execute();
    if (result.success) {
      this.commandHistory.push(command);
      if (this.commandHistory.length > this.MAX_HISTORY) {
        this.commandHistory.shift();
      }
    }
    return result;
  }

  async undoLastCommand(): Promise<CommandResult | null> {
    const lastCommand = this.commandHistory.pop();
    if (!lastCommand) return null;
    return await lastCommand.undo();
  }
}
\end{lstlisting}

\subsection{Mẫu thiết kế Rule Engine trong xử lý PCCC và tự động ngắt điện}
\label{subsec:rule_engine_pccc}

Module an toàn PCCC (\texttt{EmergencyService}) được xây dựng theo mẫu thiết kế \textbf{Rule Engine} nhằm phân loại chính xác các mức độ rủi ro và kích hoạt chuỗi hành động bảo vệ tương ứng theo thời gian thực:

\begin{itemize}
    \item \textbf{Mức NORMAL (Bình thường):} Nồng độ gas $< 50\%$ và nhiệt độ $< 40^\circ\text{C}$. Hệ thống duy trì chế độ giám sát vi khí hậu tiêu chuẩn.
    \item \textbf{Mức WARNING (Cảnh báo sớm):} Nồng độ gas trong khoảng $50\% - 79\%$ hoặc nhiệt độ trong khoảng $40^\circ\text{C} - 55^\circ\text{C}$. Hệ thống kích hoạt còi bíp cảnh báo tại chỗ, bật đèn vàng trên giao diện web và gửi thông báo cảnh giác đến sinh viên trong phòng.
    \item \textbf{Mức CRITICAL (Nguy hiểm hỏa hoạn / Rò rỉ khí gas nghiêm trọng):} Khi nồng độ gas $\ge 80\%$ hoặc nhiệt độ $\ge 60^\circ\text{C}$. Rule Engine lập tức kích hoạt chuỗi phản ứng bảo vệ nghiêm ngặt:
    \begin{enumerate}
        \item Kích hoạt còi báo động tại chỗ hú liên tục và chớp đèn LED đỏ cảnh báo.
        \item \textbf{Tự động cắt toàn bộ nguồn điện phòng:} Gửi tín hiệu ngắt toàn bộ các relay đèn và thiết bị sinh hoạt để triệt tiêu hoàn toàn nguy cơ phóng điện phát sinh tia lửa điện gây cháy nổ khí gas.
        \item Tự động bật relay quạt thông gió xả khí độc ra ngoài.
        \item Ghi nhận bản ghi sự cố khẩn cấp mức \texttt{CRITICAL} vào bảng \texttt{emergencies}.
        \item Phát sóng bản tin cảnh báo khẩn cấp tức thời (Emergency Broadcast) qua WebSockets tới toàn bộ giao diện quản trị, đồng thời gửi tin nhắn thông báo tự động qua Telegram Bot và Email SMTP đến ban quản lý KTX.
    \end{enumerate}
\end{itemize}

\section{Cơ chế xác thực và phân quyền RBAC kết hợp JWT}
\label{sec:rbac_jwt}

Hệ thống bảo đảm an toàn dữ liệu và phân tách trách nhiệm thông qua mô hình phân quyền dựa trên vai trò (Role-Based Access Control - RBAC) kết hợp chuẩn JSON Web Token (JWT):
\begin{itemize}
    \item \textbf{Vai trò ADMIN (Ban quản lý):} Có toàn quyền giám sát tất cả 18 phòng thuộc 3 tòa nhà A, B, C; quản lý danh sách sinh viên; gán mã thẻ RFID; xử lý sự cố khẩn cấp PCCC; xem toàn bộ nhật ký an ninh và xuất báo cáo thống kê.
    \item \textbf{Vai trò STUDENT (Sinh viên nội trú):} Khi đăng nhập thành công, hệ thống tự động nhận diện mã định danh phòng ở (\texttt{roomId}) của sinh viên. Sinh viên chỉ được phép xem thông số môi trường của phòng mình, điều khiển thiết bị phòng mình, xem danh sách bạn cùng phòng và xem lịch sử quẹt thẻ cá nhân. Mọi hành vi cố tình truy xuất dữ liệu phòng khác đều bị ngăn chặn ở cấp độ bảo vệ \texttt{RoomOwnerGuard}.
\end{itemize}

\section{Hệ thống thông báo đa kênh tự động}
\label{sec:thong_bao_da_kanh}

\subsection{Dịch vụ Telegram Bot Polling tương tác khẩn cấp}
Hệ thống tích hợp Telegram Bot thông qua thư viện \texttt{node-telegram-bot-api}. Khi Rule Engine phát hiện sự cố PCCC khẩn cấp, một bản tin cảnh báo với định dạng Markdown chuyên nghiệp được gửi ngay lập tức tới nhóm chỉ huy của ban quản lý KTX:

\begin{figure}[H]
\centering
\begin{tikzpicture}
    \node[draw=red!80!black, fill=red!5, rounded corners=6pt, text width=13cm, inner sep=12pt] {
        \textbf{\large [CẢNH BÁO PCCC KHẨN CẤP] — SMARTDORM TDTU}\\[0.2cm]
        \textbf{Loại sự cố:} RÒ RỈ KHÍ GAS NGUY HIỂM (\texttt{GAS\_LEAK})\\
        \textbf{Mức độ:} \textcolor{red}{\textbf{CRITICAL (CỰC KỲ NGUY CẤP)}}\\
        \textbf{Vị trí:} Phòng 101 – Tầng 1 – Tòa nhà A\\
        \textbf{Nồng độ khí gas đo được:} \textbf{100\% (Vượt ngưỡng an toàn 80\%)}\\
        \textbf{Hành động tự động đã thực hiện:} Đã tự động cắt toàn bộ nguồn điện phòng và kích hoạt quạt thông gió xả khí.\\
        \textbf{Thời gian phát hiện:} 21/10/2026 – 23:55:05\\
        \rule{\linewidth}{0.5pt}\\[0.1cm]
        \textit{[ Nút bấm tương tác trên Telegram: Xác Nhận Đã Điều Phối Đội PCCC ]}
    };
\end{tikzpicture}
\caption{Cấu trúc bản tin cảnh báo PCCC khẩn cấp tự động gửi qua Telegram Bot}
\label{fig:telegram_alert_sample}
\end{figure}

\subsection{Dịch vụ Email SMTP tự động gửi báo cáo và sự cố}
Sử dụng thư viện Nodemailer kết hợp giao thức SMTP bảo mật (TLS Port 587), hệ thống tự động biên soạn và gửi email thông báo định dạng HTML tiêu chuẩn tới địa chỉ thư điện tử của ban quản lý. Email chứa đầy đủ thông tin sự cố, giá trị cảm biến đo đạc, danh sách sinh viên đang lưu trú trong phòng gặp nạn để ban quản lý nhanh chóng tổ chức phương án hỗ trợ di tản khi cần thiết.

\newpage
